% TOP_PROBLEM: {"id":30007755,"problem_number":"LOCAL-30007755","title":"Inheritance taxation and gifts","category_id":21,"category":{"id":21,"name":"economics","display_name":"Economics","description":"Open economic theory statements, published research agendas, and proposed scoped specifications; question provenance and historical priority are qualified per record.","slug":"economics","order_index":21,"created_at":"2026-10-08T00:00:00Z"},"status":"research_question","external_url":null,"legacy_ids":[],"published":false,"statement":"In the explicitly specified setting: How can inheritance taxes reduce inherited advantage while accounting for saving responses, lifetime gifts, and avoidance? The mathematical statement above fixes the target, assumptions and scope of this catalogue formulation.","economics":{"original_id":244,"field":"Taxation and social insurance","kind":"design","formalization_basis":"adapted_research_question","source_status":"research_agenda","priority_status":"earliest_found","priority_note":"Earliest explicit agenda verified in this audit, not a claim of first historical publication. The mathematical scope is authored for this catalogue.","earliest_verified_reference":{"authors":["Arun Advani","David Sturrock"],"title":"Reforming inheritance tax","year":2023,"url":"https://ifs.org.uk/publications/reforming-inheritance-tax","doi":"10.1920/re.ifs.2023.0275","locator":"Box 7.5, Responses to tax changes; section 7.6, Effects on social mobility","evidence_summary":"Explains why static reform costings omit behavioral changes and why saving responses are difficult to estimate. This motivates a scoped empirical/design agenda; first historical priority is not established."},"current_reference":{"authors":["Arun Advani","David Sturrock"],"title":"Reforming inheritance tax","year":2023,"url":"https://ifs.org.uk/publications/reforming-inheritance-tax","doi":"10.1920/re.ifs.2023.0275","locator":"Box 7.5, Responses to tax changes; section 7.6, Effects on social mobility","evidence_summary":"Explains why static reform costings omit behavioral changes and why saving responses are difficult to estimate. This motivates a scoped empirical/design agenda; first historical priority is not established."},"formalization_note":"Catalogue-authored scoped formalization. Population, instruments, welfare objective and identification restrictions must be instantiated before claiming a resolution; source authors are not credited with these added assumptions."},"statusQualification":"Published question or research agenda verified at the cited locator. The mathematical specification is adapted for this catalogue and includes additional assumptions; source publication does not establish first-ever priority or certify this exact model remains unsolved. Catalogue-authored scoped formalization. Population, instruments, welfare objective and identification restrictions must be instantiated before claiming a resolution; source authors are not credited with these added assumptions.","statusReviewedAt":"2026-10-08"}
% FIRST_POSTED: 2026-10-08
% ENTRY_KIND: statement_only
% !TeX program = lualatex
\documentclass[11pt]{article}
\usepackage{fontspec}
\usepackage[a4paper,margin=25mm]{geometry}
\usepackage{amsmath,amssymb,mathtools}
\usepackage{xurl}
\usepackage[hidelinks]{hyperref}
\newcommand{\sref}[2]{\hyperlink{source:#1:#2}{[S#2]}}
\setlength{\emergencystretch}{3em}
\begin{document}
% problemId:30007755
\section*{Inheritance taxation and gifts}\label{30007755}
\subsection{Definitions and mathematical statement}
Fix a donor–recipient population and a finite horizon \(H\). A transfer-tax policy \(p\) specifies estate and lifetime-gift schedules, exemptions, information reporting and enforcement. For each family \(i\), let \(S_i(p)\) be donor saving, \(G_i(p)\) total lifetime gifts, \(B_i(p)\) the estate at death, and \(A_i(p)\) unreported transfers. Recipient terminal resources \(Y_i^H(p)\) include net transfers and subsequent earnings and investment. Let \(V_i(p)\) be a specified family welfare index and \(R(p)\) net public receipts. Define inherited advantage as the population regression slope
\[\beta(p)=\frac{\operatorname{Cov}(Y_i^H(p),P_i)}{\operatorname{Var}(P_i)},\]
where \(P_i\) is a fixed pre-policy parental-resources measure with positive variance. For a specified baseline \(p_0\), welfare-loss tolerance \(\delta\geq0\), funding requirement \(\bar R\), and admissible policies \(\mathcal P\), characterize policies minimizing \(\beta(p)\) subject to \(R(p)\geq\bar R\) and \(\mathbb E[V_i(p)-V_i(p_0)]\geq-\delta\). Identify or bound the required counterfactuals from linked family and transfer records, allowing different bequest motives and endogenous gift timing. State assumptions on concealed transfers and family substitution. Comparing static estate distributions alone does not resolve this policy problem.

\par\textbf{Formulation and scope.} Catalogue-authored scoped formalization. Population, instruments, welfare objective and identification restrictions must be instantiated before claiming a resolution; source authors are not credited with these added assumptions.

\par\textbf{Open-question provenance.} An explicit question or research agenda was verified at the source locator. This does not by itself certify that every later version remains unresolved \sref{30007755}{1}.

\par\textbf{Historical priority.} Earliest explicit agenda verified in this audit, not a claim of first historical publication. The mathematical scope is authored for this catalogue.

\subsection{Short English statement}
In the explicitly specified setting: How can inheritance taxes reduce inherited advantage while accounting for saving responses, lifetime gifts, and avoidance? The mathematical statement above fixes the target, assumptions and scope of this catalogue formulation.

\subsection{Sources}
\begin{itemize}\raggedright
\item[\textnormal{[S1]}]\hypertarget{source:30007755:1}{} Arun Advani, David Sturrock. 2023. Reforming inheritance tax. Box 7.5, Responses to tax changes; section 7.6, Effects on social mobility.
\url{https://ifs.org.uk/publications/reforming-inheritance-tax}.
DOI: 10.1920/re.ifs.2023.0275.
{\small\textit{Earliest verified question/agenda reference; historical priority qualified below. Explains why static reform costings omit behavioral changes and why saving responses are difficult to estimate. This motivates a scoped empirical/design agenda; first historical priority is not established.}}
\end{itemize}
\end{document}
